\section{从线性分类器到神经网络}

\subsection{线性分类器的局限性}

线性分类器 $f = W\mathbf{x}$ 只能学习一条直线（或超平面）作为决策边界。对于线性不可分的数据（如环形分布），线性分类器完全无法处理。

\begin{figure}[H]
\centering
\includegraphics[width=0.95\textwidth]{slides-images/slide-023.jpg}
\caption{线性不可分问题：左侧的环形数据无法用一条直线分开\protect\footnotemark}
\end{figure}
\footnotetext{Slides 第24页。}

\begin{importantbox}{非线性变换的力量}
解决方案是将数据从原始空间\textbf{非线性变换}到新空间。例如，将笛卡尔坐标 $(x, y)$ 转换为极坐标 $(r, \theta)$ 后，原本线性不可分的环形数据变得线性可分。神经网络的核心能力就是\textbf{自动学习}这种非线性变换。
\end{importantbox}

\subsection{构建两层神经网络}

最简单的神经网络——两层全连接网络的公式为：

$$f = W_2 \cdot \max(0, W_1 \mathbf{x})$$

\begin{figure}[H]
\centering
\includegraphics[width=0.95\textwidth]{slides-images/slide-021.jpg}
\caption{从一层到两层：引入隐藏层和非线性激活函数\protect\footnotemark}
\end{figure}
\footnotetext{Slides 第22页。}

关键维度分析：
\begin{itemize}
    \item $\mathbf{x} \in \mathbb{R}^D$：输入向量（$D$ 个特征）
    \item $W_1 \in \mathbb{R}^{H \times D}$：第一层权重（$H$ 个隐藏神经元）
    \item $W_2 \in \mathbb{R}^{C \times H}$：第二层权重（$C$ 个输出类别）
    \item $\max(0, \cdot)$：ReLU 激活函数（非线性）
\end{itemize}

\begin{warningbox}{没有非线性等于没有多层}
如果去掉 $\max(0, \cdot)$，两层网络变成 $f = W_2 \cdot W_1 \mathbf{x}$，而 $W_2 W_1$ 可以合并为一个矩阵 $W_3$，回到单层线性分类器！非线性激活函数是神经网络能够解决非线性问题的\textbf{根本原因}。
\end{warningbox}

\subsection{更深的网络}

可以继续堆叠更多层来构建更深的网络：

$$f = W_3 \cdot \max(0, W_2 \cdot \max(0, W_1 \mathbf{x}))$$

\begin{figure}[H]
\centering
\includegraphics[width=0.95\textwidth]{slides-images/slide-022.jpg}
\caption{三层神经网络：每层之间插入非线性激活函数\protect\footnotemark}
\end{figure}
\footnotetext{Slides 第23页。}

这些只使用全连接层（矩阵乘法）和激活函数的网络被称为\textbf{全连接网络}（Fully Connected Network）或\textbf{多层感知机}（MLP, Multi-Layer Perceptron）。

\subsection{模板学习的视角}

\begin{figure}[H]
\centering
\includegraphics[width=0.95\textwidth]{slides-images/slide-025.jpg}
\caption{模板视角：线性分类器每类学一个模板，神经网络可以学多个部分模板\protect\footnotemark}
\end{figure}
\footnotetext{Slides 第26页。}

从模板的角度理解网络能力的提升：

\begin{itemize}
    \item \textbf{线性分类器}：每个类别只有 1 个模板（即 $W$ 的对应行），很受限
    \item \textbf{神经网络}：隐藏层有 $H$ 个神经元，可以学到 $H$ 个``部分模板''。例如，鸟、猫、鹿、狗都有眼睛，某个隐藏神经元可能专门检测``眼睛''特征，被多个类别共享
\end{itemize}

\begin{importantbox}{神经网络的表示能力}
增加隐藏层神经元的数量相当于增加网络的``模板库''。更多神经元 = 更多部分模板 = 更强的函数逼近能力 = 更复杂的决策边界。但过多的容量也会导致过拟合——这时需要正则化来控制。
\end{importantbox}

\subsection{本章小结}

神经网络通过交替使用线性变换（矩阵乘法）和非线性激活函数来突破线性分类器的局限。两层网络 $f = W_2 \cdot \sigma(W_1 \mathbf{x})$ 已经可以理论上逼近任意连续函数。隐藏层的宽度决定了网络的表示能力。

%% ========== 激活函数 ==========

